\chapter{Research Questions, Limitations and Conclusions}

The preceding analyses can now be combined to answer the research questions
introduced by the project. The objective of this chapter is not to introduce
additional models, but to synthesize evidence from the EDA and the different
statistical-learning approaches.

% -----------------------------------------------------------------------------
\section{RQ1: Are Arrest Outcomes Distributed Uniformly Across Recorded Stops?}

\textbf{Answer: No.} The descriptive analysis provides clear evidence that
arrest outcomes are heterogeneous across demographic, temporal, and
geographical characteristics.

Among two-wheeled stops, the observed arrest rate is $12.1\%$ in the Bronx
but only approximately $3.2\%$ in Manhattan. For four-wheeled stops the
geographical range is narrower but remains substantial, from approximately
$4.8\%$ in the Bronx to $2.5\%$ in Queens.

Time of day produces an especially consistent pattern. At approximately
07:00, arrest rates fall to $3.2\%$ for two-wheeled vehicles and $0.6\%$
for four-wheeled vehicles. By 03:00, they rise to approximately $10.9\%$
and $5.3\%$, respectively.

Demographic differences are also substantial. For two-wheeled vehicles,
Hispanic and Black motorists have observed arrest rates of $7.24\%$ and
$6.86\%$, compared with $3.23\%$ for White motorists. For four-wheeled
vehicles, the corresponding rates are $4.66\%$, $5.02\%$, and $1.15\%$.

The EDA therefore rejects the hypothesis that arrest outcomes are uniformly
distributed across the observed characteristics. These are, however,
\textbf{unadjusted differences}; they do not establish which characteristics
retain an association after accounting for the others.

% -----------------------------------------------------------------------------
\section{RQ2: Do Demographic Disparities Persist After Adjustment?}

\textbf{Answer: Yes.} Several of the demographic differences identified by the
EDA remain after conditioning on temporal, geographical, and other demographic
characteristics.

The clearest evidence comes from reported ethnicity. In the Logistic
Regression for two-wheeled vehicles,

$$
OR_{\text{Black}}=1.451,
\qquad
OR_{\text{Hispanic}}=1.778,
$$

relative to the White reference category.

For four-wheeled vehicles the estimated associations are much larger:

$$
OR_{\text{Black}}=3.425,
\qquad
OR_{\text{Hispanic}}=3.258.
$$

These findings are not dependent on a single model. In the four-wheeled LASSO
1-SE model, the corresponding penalized odds ratios remain approximately

$$
3.202
\quad\text{and}\quad
3.057,
$$

while the GAM estimates

$$
3.349
\quad\text{and}\quad
3.185.
$$

The same consistency is visible for sex. Female observations have conditional
odds ratios of $0.300$ and $0.408$ in the two Logistic models, and the
direction and magnitude remain similar under LASSO and GAM.

Therefore, demographic disparities do not disappear after controlling for the
observed context included in the models. In particular, the Black and Hispanic
coefficients remain large and resistant to regularization in the four-wheeled
analysis.

These results demonstrate \textbf{persistent conditional statistical
disparities}. They do not, by themselves, identify the mechanism responsible
for those disparities.

% -----------------------------------------------------------------------------
\section{RQ3: Are the Relationships with Hour and Age Non-Linear?}

\textbf{Answer: Yes.} The GAM provides strong statistical and graphical
evidence that both relationships are non-linear.

For hour of day,

$$
\text{edf}_{2W}=7.52,
\qquad
\text{edf}_{4W}=7.925,
$$

with $p<2\times10^{-16}$ in both models. The smooth functions reproduce the
same broad cycle observed during the EDA: elevated arrest propensity during
late-night and early-morning hours, a pronounced minimum around 07:00--08:00,
and a progressive increase during the afternoon and evening.

The consistency of this curve across vehicle types makes time of day one of
the most stable findings in the entire analysis.

Age is also non-linear,

$$
\text{edf}_{2W}=4.66,
\qquad
\text{edf}_{4W}=4.878,
$$

again with $p<2\times10^{-16}$.

Unlike hour, however, the age curves differ substantially between vehicle
types. The four-wheeled smooth reaches its highest value around ages 27--30
and subsequently declines, particularly after age 50. The two-wheeled curve
has a more irregular structure, with elevated values among the youngest
observations and another local increase at later ages.

The hypothesis that age and hour can be represented adequately by simple
linear effects is therefore rejected.

% -----------------------------------------------------------------------------
\section{RQ4: Do Two-Wheeled and Four-Wheeled Stops Exhibit the Same Patterns?}

\textbf{Answer: Only partially.}

Some results are remarkably consistent. Both datasets exhibit:

\begin{itemize}
\item a strong late-night versus morning arrest pattern;
\item lower conditional odds for female motorists;
\item positive conditional associations for Black and Hispanic categories
relative to White;
\item substantial geographical heterogeneity;
\item only moderate individual-level predictability.
\end{itemize}

However, the magnitude and sometimes even the direction of specific
associations differ considerably.

The ethnicity contrast is substantially stronger among four-wheeled vehicles.
For example,

$$
OR_{\text{Black}}:
1.451\;(2W)
\quad\text{vs.}\quad
3.425\;(4W),
$$

and

$$
OR_{\text{Hispanic}}:
1.778\;(2W)
\quad\text{vs.}\quad
3.258\;(4W).
$$

Age provides an even clearer difference. In the Logistic Regression, Under-18
two-wheeled observations have higher estimated odds than the 25--34 reference
group,

$$
OR=1.725,
$$

whereas the corresponding four-wheeled estimate is

$$
OR=0.756.
$$

Similarly, the four-wheeled model exhibits a strong progressive reduction at
older ages, reaching $OR_{65+}=0.201$, while the corresponding two-wheeled
65+ coefficient is not statistically distinguishable from the 25--34
reference group.

The GAM smooths confirm that this is not merely an artifact of the age
brackets.

Consequently, the data support the decision to maintain separate analyses.
There appears to be a common temporal component, but the demographic
structure of arrest outcomes is strongly dependent on the vehicle context.

% -----------------------------------------------------------------------------
\section{RQ5: How Predictable Is an Arrest from the Available Information?}

\textbf{Answer: Arrest is moderately predictable as a ranking problem, but not
with high reliability at the individual classification level. Predictability
is consistently greater for four-wheeled stops.}

Across two-wheeled models, ROC-AUC ranges approximately from

$$
0.682 \text{ to } 0.690,
$$

whereas for four-wheeled models it remains around

$$
0.732 \text{ to } 0.735.
$$

Thus, all models perform better than random ranking, but none achieves strong
class separation.

The PR-AUC results must be interpreted relative to the different class
prevalences. For two-wheeled vehicles, PR-AUC is approximately
$0.117$--$0.125$ against a no-skill baseline of approximately $0.063$.
For four-wheeled vehicles it is approximately $0.076$--$0.079$ against a
much lower baseline of approximately $0.034$.
odels. Consequently, most observations
classified as arrests are false positives even though sensitivity is generally
around $70\%-76\%$.

The models therefore contain genuine information about relative arrest
propensity, but the observed characteristics are not sufficient to determine
individual arrest outcomes with high certainty.

The GAM fit reinforces this conclusion from another perspective: only about
$9.61\%$ and $9.31\%$ of deviance is explained in the two vehicle
populations. A substantial part of the outcome therefore depends on
information that is either unobserved or not represented by the available
predictors.

% -----------------------------------------------------------------------------
\section{RQ6: Do the Results Demonstrate Systematic or Excessive Enforcement?}

\textbf{Answer: The results provide evidence of persistent systematic
statistical disparities in post-stop arrest outcomes, but they are not
sufficient to establish systematic discriminatory or excessive enforcement.}

There is substantial empirical evidence that the observed differences are not
purely descriptive artifacts. For example, among four-wheeled vehicle stops,
the raw arrest rate is $5.02\%$ for Black motorists and $4.66\%$ for
Hispanic motorists compared with $1.15\%$ for White motorists. After
adjustment, the Logistic Regression still estimates odds ratios of $3.425$
and $3.258$, respectively. LASSO and GAM produce very similar conditional
relationships.

Similarly, significant geographical, sex-related, age-related, and temporal
differences remain after adjustment, and several of these patterns are
recovered under multiple modelling assumptions.

This constitutes evidence that the disparities are
\textbf{persistent and structured within the recorded stop data}.

However, three different statements must be distinguished:

$$
\boxed{
\text{Observed disparity}
\neq
\text{causal demographic effect}
\neq
\text{proof of discriminatory enforcement}
}
$$

The dataset contains only vehicle stops that have already occurred. It
therefore allows the analysis of

$$
P(\text{arrest}\mid\text{recorded stop},X),
$$

but not the probability that a motorist is stopped in the first place.

Furthermore, the records do not contain complete information about the reason
for the stop, the legal circumstances leading to the arrest, the behavior
during the encounter, or other operational factors that may influence the
decision to arrest. Consequently, important confounding information is not
available to the models.

The strongest conclusion supported by the available evidence is therefore the
following:

\begin{quote}
Among recorded NYPD vehicle stops, arrest outcomes exhibit persistent
demographic, temporal, and geographical disparities. Several of these
differences remain large after adjustment for the observed characteristics and
are reproduced across alternative statistical models, particularly in the
four-wheeled population. However, the available data are insufficient to
determine whether these conditional disparities are caused by discriminatory
or excessive police enforcement.
\end{quote}

This distinction is central to the interpretation of the project. The analysis
does not confirm a predetermined claim of police bias, nor does it provide
evidence that the observed disparities are innocuous. Instead, it identifies
which disparities are present, quantifies their magnitude, tests their
robustness across modelling assumptions, and establishes the boundary beyond
which the available data cannot support stronger causal conclusions.

% =============================================================================
% FINAL SYNTHESIS
% =============================================================================

% -----------------------------------------------------------------------------
\section{Limitations}

TODO

% -----------------------------------------------------------------------------
\section{Synthesis and Conclusions}

Taken together, the exploratory and modelling analyses produce a coherent
picture.

The most stable finding is temporal: arrest propensity follows a pronounced
daily cycle in both vehicle populations, with substantially lower conditional
odds during morning hours and higher values during the night. The agreement
between raw arrest rates, Logistic Regression coefficients, LASSO retention,
Naive Bayes conditional distributions, and GAM smooths makes this the most
robust cross-dataset pattern.

Demographic variables also contain persistent information, but their effects
are considerably more context-dependent. Reported ethnicity and sex remain
associated with arrest after adjustment and regularization, while the magnitude
of the ethnicity coefficients is substantially greater among four-wheeled
vehicle stops. Age exhibits an even stronger vehicle-specific structure.

At the same time, predictive performance places an important limit on the
interpretation. The available predictors produce moderate ranking ability but
low individual precision, and more flexible modelling does not substantially
raise this performance ceiling.

The principal contribution of the modelling stage is therefore not the
identification of a single ``best'' classifier. It is the convergence of
several modelling approaches toward the same broader conclusion:
\textbf{post-stop arrest outcomes exhibit structured and persistent
disparities, but much of the individual outcome remains unexplained by the
information available in the dataset.}

TODO